\section{A uniform functional law for the mechanical clock}
\label{sec:uniform-fluid-clock}
Finite-record continuity also has a long-time consequence that does not
require a mixing rate. We prove a parameter-uniform functional law at the
first-order scale, including the maximum unfinished return over an entire
observation interval. The \(\sqrt t\)-scale process theorem in
Proposition~\ref{prop:clock} remains a separate statement.

Set
\[
 G_R=(\kappa_R^*,r_R^*,\tau_R^*),\qquad
 \bar G_R=(0,0,c_*^{-1},\bar\tau_R/c_*),\qquad J_{0,R}=0.
\]
All probabilities on the return section use $\nu_R^*$, not an
independent sequence of records. Write $\bar r=c_*^{-1}$ and
$\bar\tau_R^*=\bar\tau_R/c_*$. Both $\bar\tau_R^*$ and its reciprocal
are bounded on $I$.

\begin{lemma}[A compact subadditive argument]
\label{lem:compact-subadditive}
Let $a_n:I\to[0,\infty)$ be continuous, $a_0=0$, and
$a_{n+m}(R)\le a_n(R)+a_m(R)$. Suppose
$\inf_{n\ge1}a_n(R)/n=0$ for every $R$. Then
\[
 \lim_{n\to\infty}\sup_{R\in I}\frac{a_n(R)}n=0.
\]
More precisely, for every $\epsilon>0$ there is a finite integer $M$
such that, with $C=\sup_R a_1(R)$,
\begin{equation}\label{eq:subadditive-uniform-budget}
 \sup_R a_n(R)/n\le\epsilon+CM/n\qquad(n\ge1).
\end{equation}
\end{lemma}
\begin{proof}
For each $R$ choose a block length $m_R$ with
$a_{m_R}(R)/m_R<\epsilon$. Continuity preserves this strict inequality
on a neighborhood of $R$. Take a finite subcover and let $M$ be the
largest of its block lengths. For an arbitrary parameter choose one of
the covering neighborhoods, with block length $m\le M$, and write
$n=qm+r$, $0\le r<m$. Subadditivity and $a_r\le r a_1$ give
$a_n\le q\epsilon m+Cr\le\epsilon n+CM$. This proves the
estimate and then the limit. The neighborhoods need not have a common
block length; replacing them by an unproved common mixing time would
not be justified.
\end{proof}

\begin{theorem}[Uniform first-order return process]
\label{thm:uniform-return-fluid}
The actual induced records satisfy
\begin{equation}\label{eq:uniform-return-mean-law}
 \lim_{n\to\infty}\sup_{R\in I}
       \int\left|J_{n,R}/n-\bar G_R\right|\dd\nu_R^*=0.
\end{equation}
For every fixed $A<\infty$ and $\epsilon>0$,
\begin{equation}\label{eq:uniform-return-functional}
 \lim_{n\to\infty}\sup_{R\in I}\nu_R^*\left\{
 \sup_{0\le u\le A}
 \left|n^{-1}J_{\lfloor nu\rfloor,R}-u\bar G_R\right|>\epsilon
 \right\}=0.
\end{equation}
\end{theorem}
\begin{proof}
The ergodic collision map in Section~\ref{sec:return-stability}
induces an ergodic map on its positive-measure return section. Indeed,
the complete tower over an invariant subset of the induced map is an
invariant collision set. The two towers over that subset and its
complement partition the full collision space modulo null sets;
collision ergodicity forces one of the two to have zero measure.
Thus the integrable vector $G_R$ obeys the $L^1$ ergodic theorem for
each fixed $R$, with the exact mean in
Proposition~\ref{prop:exact-induced-means}.

Put $a_n(R)=\int|J_{n,R}-n\bar G_R|\dd\nu_R^*$.
Invariance of $F_R^*$ and the cocycle identity make $a_n$ subadditive.
It is continuous: on the common collision space it equals
\[
 c_*^{-1}\int_M
 \left|\widetilde J_{n,R}-n\bar G_R\one_{Y_R^*}\right|\dd\nu.
\]
Theorem~\ref{thm:return-stability} gives the first term's $L^1$
continuity, the section indicators are continuous in $L^1$ because
their boundaries move continuously and have zero measure, and the
mean is continuous by its explicit formula. The pointwise ergodic
theorem gives $a_n(R)/n\to0$. Lemma~\ref{lem:compact-subadditive}
proves \eqref{eq:uniform-return-mean-law}.

For the functional conclusion, first test a fixed finite mesh of
$[0,A]$. Equation~\eqref{eq:uniform-return-mean-law} and Markov's
inequality give convergence in probability at all its mesh points,
uniformly in $R$. The count and roof coordinates of $J_n$ are
nondecreasing. Between mesh points their normalized deviations are
therefore bounded by the endpoint deviations plus a uniformly bounded
mean times the mesh width, with an $O(n^{-1})$ integer-part error.
For the displacement coordinate, finite horizon gives a constant
$C_\kappa$ such that for every $l\ge j$,
\[
 |K_{l,R}-K_{j,R}|\le C_\kappa(N_{l,R}-N_{j,R}).
\]
Consequently its oscillation on a mesh interval is controlled by the
count increment on that interval. Its mean is zero. On the event
where all mesh errors are small, every coordinate has a uniformly
small error between mesh points as well. Let $n$ grow, then let the
mesh width and the mesh error tolerance decrease. This proves
\eqref{eq:uniform-return-functional} without an independence or
maximal-martingale assumption.
\end{proof}

\begin{lemma}[Sublinear maximum return blocks]
\label{lem:uniform-maximal-block}
Put $Q_R=r_R^*+\tau_R^*+|\kappa_R^*|$ and
\[
 U(L)=\sup_R\int Q_R\one_{\{Q_R>L\}}\dd\nu_R^*.
\]
Then $U(L)\to0$. For $n\ge1$, $A\ge0$, and $\epsilon>0$,
\begin{equation}\label{eq:uniform-maximal-block}
 \sup_R\nu_R^*\left\{
 \max_{0\le j\le\lfloor An\rfloor}Q_R\circ(F_R^*)^j>\epsilon n
 \right\}\le\frac{A+1}{\epsilon}U(\epsilon n).
\end{equation}
Both \eqref{eq:uniform-return-functional} and the convergence to zero
in \eqref{eq:uniform-maximal-block} remain valid when the initial
base point has the stationary length-biased law
\begin{equation}\label{eq:length-biased-base-v4}
 \dd\rho_R=\frac{\tau_R^*}{\bar\tau_R^*}\dd\nu_R^*.
\end{equation}
\end{lemma}
\begin{proof}
Finite horizon bounds $\tau_R^*$ and $|\kappa_R^*|$ by a uniform
constant times $r_R^*$. The uniform integrability in
Theorem~\ref{thm:return-stability}, and the constant section mass,
therefore give $U(L)\to0$. Invariance of $F_R^*$, the union bound
and truncated Markov inequality give
\[
 \nu_R^*\{\max_{j\le\lfloor An\rfloor}Q_R\circ F^j>\epsilon n\}
 \le (\lfloor An\rfloor+1)\nu_R^*(Q_R>\epsilon n)
 \le\frac{A+1}{\epsilon}\int Q_R\one_{Q_R>\epsilon n}\dd\nu_R^*.
\]
No independence of the events is used.

Let $d_*:=\inf_R\bar\tau_R^*>0$. For any event $B_R$ determined
by the forward base orbit and every $L>0$,
\begin{equation}\label{eq:length-bias-transfer-v4}
 \rho_R(B_R)\le d_*^{-1}\left(
 L\nu_R^*(B_R)+\int\tau_R^*\one_{\{\tau_R^*>L\}}\dd\nu_R^*\right).
\end{equation}
If $\sup_R\nu_R^*(B_R)\to0$, first hold $L$ fixed and pass to that
limit, then let $L\to\infty$. Uniform integrability of the roof
makes the right side tend to zero uniformly. Apply this argument to
the exceptional events in the two displayed convergence statements.
This transfer uses a length-biased law, not the unweighted launch law.
\end{proof}

For the stationary physical flow, let $V_R(t)$ count future visits to
$Y_R^*$, let $C_R(t)$ count future physical collisions, and let
$W_R(t)\in\Z^2$ be the net fundamental-cell displacement since time
zero. Changing the choice of a fixed fundamental cell changes the
continuous-position comparison only by a bounded term.

\begin{theorem}[Uniform stationary physical clock]
\label{thm:uniform-physical-clock}
For every $\epsilon>0$, under the actual stationary billiard laws
$\mu_R$, uniformly for $R\in I$,
\begin{equation}\label{eq:uniform-physical-clock}
 \sup_{0\le u\le1}\left|
 \frac1t\bigl(V_R(tu),C_R(tu),W_R(tu)\bigr)
 -\left(\frac{uc_*}{\bar\tau_R},
        \frac{u}{\bar\tau_R},0\right)\right|
 \longrightarrow0
 \quad\hbox{in probability as }t\to\infty.
\end{equation}
The maximum marked unfinished return encountered during $[0,t]$,
including the initial return containing time zero, is $o_P(t)$
uniformly in $R$.
\end{theorem}
\begin{proof}
Represent the starting state by $(x,s)$ in the induced suspension,
$0\le s<\tau_R^*(x)$. The marginal law of $x$ is exactly $\rho_R$
in \eqref{eq:length-biased-base-v4}. Write $T_{j,R}$ for the roof
coordinate of $J_{j,R}$. Count future events in $(0,t]$. With this
convention the following inverse identity also holds at event times:
\[
 V_R(tu)=\max\{j\ge0:T_{j,R}(x)\le s+tu\}.
\]
Choose a fixed $A$ with $A d_*>2$. The functional return law,
transferred by Lemma~\ref{lem:uniform-maximal-block}, shows that
$T_{\lfloor At\rfloor,R}>s+t$ with probability tending to one
uniformly. Here $s/t\to0$ uniformly because $s\le Q_R(x)$ and
the same lemma applies at $j=0$. Integer parts do not affect these
statements; the return-process theorem may be applied with
$n=\lceil t\rceil$ and a slightly larger fixed horizon.

On this event all inverse indices lie below $At$. Uniformly for
$0\le j\le At$ we have
$|T_{j,R}-j\bar\tau_R^*|/t=o_P(1)$ and
$\max_{j\le At}Q_R\circ F^j/t=o_P(1)$. The inequalities
\[
 T_{V_R(tu),R}\le s+tu<T_{V_R(tu)+1,R}
\]
therefore imply
\[
 \sup_{u\le1}\left|
 \bar\tau_R^* V_R(tu)/t-u\right|=o_P(1)
\]
uniformly in $R$. Division by $\bar\tau_R^*\ge d_*$ yields the
first component in \eqref{eq:uniform-physical-clock}.

The completed-return count $N_{V_R(tu),R}$ differs from the actual
physical count by the initial prefix and the final incomplete return.
Their sum is at most a fixed constant plus twice the maximal $Q_R$
over these return indices. The displacement comparison has the same
bound, up to a uniform multiplicative constant, since each physical
lattice step is bounded. These bounds hold simultaneously for all
$u\in[0,1]$. The functional return law at the random inverse indices
now gives
\[
 C_R(tu)/t=\bar r\,V_R(tu)/t+o_P(1),\qquad
 W_R(tu)/t=o_P(1)
\]
with uniform suprema. Since $\bar r/\bar\tau_R^*=1/\bar\tau_R$,
this proves the remaining components. Every unfinished piece in
$[0,t]$ belongs to one of the same return blocks and is bounded by
its marked size, proving the final assertion.
\end{proof}

\begin{corollary}[Geometric error in the physical collision rate]
\label{cor:physical-rate-error}
Let $b_*:=\min_{R\in I}\bar\tau_R$ and put
$L_*=15/(4b_*^2)$. There is a function $\omega_t(\epsilon)\to0$
for every fixed $\epsilon>0$ such that, for any $I$-valued radius
estimate $\widehat R$ on a probability space carrying the stationary
billiard at the true radius $R$, and any $\delta>0$,
\begin{align}\label{eq:physical-rate-error}
 \Pr_R\left\{\sup_{0\le u\le1}
 \left|C_R(tu)/t-u/\bar\tau_{\widehat R}\right|
                       >\epsilon+L_*\delta\right\}
 \le\omega_t(\epsilon)+\Pr_R\{|\widehat R-R|>\delta\}.
\end{align}
No independence of $\widehat R$ and the observed orbit is required.
For the visit rate, replace the deterministic error constant by $c_*L_*$.
\end{corollary}
\begin{proof}
The explicit mean satisfies $|\bar\tau_R'|<15/4$, so
$|(1/\bar\tau_R)'|\le L_*$. On the event
$|\widehat R-R|\le\delta$ the two proposed rates differ by at
most $L_*\delta$, uniformly over $u\in[0,1]$. The triangle
inequality and Theorem~\ref{thm:uniform-physical-clock} prove the
claim with $\omega_t$ equal to the supremum, over $R$, of the true
rate's exceptional probability. Multiplication by $c_*$ gives the
visit-rate statement.
\end{proof}

The supremum statement here is at scale $t$. Uniform integrability
alone does not turn it into an $o_P(\sqrt t)$ bound for a growing
window. Likewise \eqref{eq:subadditive-uniform-budget} and
\eqref{eq:physical-rate-error} do not supply a numerical convergence
rate for $\omega_t$ without quantitative finite-block estimates.
They establish the full first-order clock comparison for this moving
family. The Gaussian covariance and processes require the additional
collision estimates proved in Sections~\ref{sec:collision-covariance}--
\ref{sec:functional-gaussian}; raw density and shrinking-conditioning
conclusions require the further local estimates stated there. The
stationary law throughout is not substituted for an unknown experimental
preparation law.
